\documentclass[../../main.tex]{subfiles}% 注意这里的文件路径不能用 ./main.tex ,否则用latexmk编译子文件会报错
\graphicspath{{\subfix{./image/}}{\subfix{../../image/}}} % 指定图片引用路径,后续可以直接使用图片文件名
% 如果图片存放在上级目录,则使用 ../../image/ 作为备用路径

% 例如:
% \begin{figure}[H]
% \centering
% \includegraphics[width=0.3\textwidth]{图.png}
% \caption{}
% \label{figure:图}
% \end{figure}
% 注意:上述\label{}一定要放在\caption{}之后,否则引用图片序号会只会显示??.

\begin{document}

\section{重积分计算}

\subsection{常规计算}

\begin{theorem}[二重积分的Newton-Leibniz公式]
设 $D = [a,b] \times [c,d] \subset \mathbb{R}^2$ 是闭矩形区域，函数 $f(x,y)$ 在 $D$ 上连续。若存在二元函数 $F(x,y)$ 在 $D$ 上具有连续混合偏导数 $\frac{\partial^2 F}{\partial x \partial y}$，且满足
\[
\frac{\partial^2 F}{\partial x \partial y}(x,y) = f(x,y), \qquad \forall (x,y) \in D,
\]
则
\[
\iint_D f(x,y) \, \mathrm{d}x \mathrm{d}y = F(b,d) - F(b,c) - F(a,d) + F(a,c).
\]
\end{theorem}

\begin{theorem}[平面曲线积分的分部积分法]
设 $L$ 是平面上的有向光滑曲线，起点为 $A$，终点为 $B$。函数 $u = u(x,y)$ 与 $v = v(x,y)$ 在 $L$ 上具有一阶连续偏导数。则
\[
\int_L u \, \mathrm{d}v = uv \Big|_A^B - \int_L v \, \mathrm{d}u,
\]
其中 $\mathrm{d}v = v_x \, \mathrm{d}x + v_y \, \mathrm{d}y$，$\mathrm{d}u = u_x \, \mathrm{d}x + u_y \, \mathrm{d}y$，且
\[
uv \Big|_A^B = u(B)v(B) - u(A)v(A).
\]
\end{theorem}

\begin{example}
设$f(x,y)$在$D$上有定义，且在某一圆盘外为$0$。求证：
\begin{align*}
\iint_D \left[ \frac{\partial^2 f}{\partial x^2} \frac{\partial^2 f}{\partial y^2} - \left( \frac{\partial^2 f}{\partial x \partial y} \right)^2 \right] \mathrm{d}x \mathrm{d}y = 0.
\end{align*}
\end{example}
\begin{proof}
由题设可知$f$在$\partial D$上恒为$0$，进而
$$\frac{\partial f}{\partial x}(x,y) = \frac{\partial f}{\partial y}(x,y) = 0, \quad \forall (x,y) \in \partial D.$$
由Green公式可得
\begin{align*}
&\iint_D \left[ \frac{\partial^2 f}{\partial x^2} \frac{\partial^2 f}{\partial y^2} - \left( \frac{\partial^2 f}{\partial x \partial y} \right)^2 \right] \mathrm{d}x \mathrm{d}y
= \iint_D (f_{xx} f_{yy} - f_{xy}^2) \mathrm{d}x \mathrm{d}y \\
&= \iint_D \left[ (f_{xx} f_y)_y - (f_{xy} f_y)_x \right] \mathrm{d}x \mathrm{d}y 
\xlongequal{\text{Green公式}} \oint_{\partial D} f_{xy} f_y \mathrm{d}x + f_{xx} f_y \mathrm{d}y = 0.
\end{align*}

\end{proof}

\begin{example}
设不全为0的实数$\alpha,\beta,\gamma$，考虑
\begin{align*}
C:
\begin{cases}
x^2 + y^2 + z^2 = 1 \\
\alpha x + \beta y + \gamma z = 0
\end{cases}.
\end{align*}
计算$\int_C y(y - x)\mathrm{d}s$.
\end{example}
\begin{solution}
记$K\triangleq \sqrt{\alpha^2+\beta^2}$，$M\triangleq \sqrt{\alpha^2+\beta^2+\gamma^2}$，考虑如下正交变换
\begin{align*}
\begin{pmatrix}
x'\\
y'\\
z'
\end{pmatrix}
=
\begin{pmatrix}
\frac{\alpha\gamma}{KM} & \frac{\beta\gamma}{KM} & -\frac{K}{M}\\
-\frac{\beta}{K} & \frac{\alpha}{K} & 0\\
\frac{\alpha}{M} & \frac{\beta}{M} & \frac{\gamma}{M}
\end{pmatrix}
\begin{pmatrix}
x\\
y\\
z
\end{pmatrix}
= T
\begin{pmatrix}
x\\
y\\
z
\end{pmatrix}.
\end{align*}
于是曲线$C$方程等价于
\begin{align*}
\begin{cases}
\left( \frac{\alpha\gamma}{KM}x'-\frac{\beta}{K}y' \right)^2+\left( \frac{\beta\gamma}{KM}x'+\frac{\alpha}{K}y' \right)^2+\left( -\frac{K}{M}x' \right)^2=1\\
z' = 0
\end{cases}
\Longleftrightarrow
\begin{cases}
x'^2+y'^2=1\\
z' = 0
\end{cases}.
\end{align*}
因为正交变换是等距变换，进而保持弧长微元，即
\begin{align*}
\mathrm{d}s=\sqrt{(\mathrm{d}x)^2+(\mathrm{d}y)^2+(\mathrm{d}z)^2}
= \left| \begin{pmatrix} \mathrm{d}x\\ \mathrm{d}y\\ \mathrm{d}z \end{pmatrix} \right|
= \left| T\begin{pmatrix} \mathrm{d}x'\\ \mathrm{d}y'\\ \mathrm{d}z' \end{pmatrix} \right|
= \left| \begin{pmatrix} \mathrm{d}x'\\ \mathrm{d}y'\\ \mathrm{d}z' \end{pmatrix} \right|
= \sqrt{(\mathrm{d}x')^2+(\mathrm{d}y')^2+(\mathrm{d}z')^2}
= \mathrm{d}s'.
\end{align*}
又由曲线$C:\begin{cases}x'^2+y'^2=1\\z'=0\end{cases}$的对称性知
\begin{align*}
\int_{x'^2+y'^2=1,z'=0}x'y'\mathrm{d}s'
= \int_{x'^2+y'^2=1,z'=0}(-x')y'\mathrm{d}s'
\Longrightarrow
\int_{x'^2+y'^2=1,z'=0}x'y'\mathrm{d}s'=0.
\end{align*}
所以
\begin{align*}
\int_C y(y-x)\mathrm{d}s
=& \int_{x'^2+y'^2=1,z'=0}\left( \frac{\beta\gamma}{KM}x'+\frac{\alpha}{K}y' \right)\left( \frac{(\beta-\alpha)\gamma}{KM}x'+\frac{\alpha+\beta}{K}y' \right)\mathrm{d}s'\\
=& \int_{x'^2+y'^2=1,z'=0}\left[ \frac{\beta(\beta-\alpha)\gamma^2}{K^2M^2}x'^2+\frac{\alpha(\alpha+\beta)}{K^2}y'^2 \right]\mathrm{d}s'\\
=& \int_0^{2\pi}\left[ \frac{\beta(\beta-\alpha)\gamma^2}{K^2M^2}\cos^2\theta+\frac{\alpha(\alpha+\beta)}{K^2}\sin^2\theta \right]\mathrm{d}\theta\\
=& \frac{\beta(\beta-\alpha)\gamma^2}{K^2M^2}\pi + \frac{\alpha(\alpha+\beta)}{K^2}\pi
= \pi \frac{\alpha^2+\alpha\beta+\gamma^2}{\alpha^2+\beta^2+\gamma^2}.
\end{align*}
\end{solution}

\begin{example}
对 $t\in(0,+\infty)$，令
\begin{align*}
I(t)=\iint_{\Sigma_t} P\mathrm{d}y\mathrm{d}z+Q\mathrm{d}z\mathrm{d}x+R\mathrm{d}x\mathrm{d}y,
\end{align*}
其中 $\Sigma_t$ 是柱体 $\{(x,y,z)\mid x^2+y^2\leqslant t^2,0\leqslant z\leqslant 1\}$ 的表面，定向为外侧.$P = Q = R = f\big((x^2+y^2)z\big), f \in C^1(\mathbb{R})$.求极限：
\begin{align*}
\lim\limits_{t\to0^+}\frac{I(t)}{t^4}.
\end{align*}
\end{example}
\begin{solution}
由\hyperref[theorem:Gauss公式--数分]{Gauss公式}可得
\begin{align*}
I(t)=\iiint_V \left(\frac{\partial P}{\partial x}+\frac{\partial Q}{\partial y}+\frac{\partial R}{\partial z}\right)\mathrm{d}x\mathrm{d}y\mathrm{d}z=\iiint_V f'\big((x^2+y^2)z\big)\big(2xz+2yz+x^2+y^2\big)\mathrm{d}x\mathrm{d}y\mathrm{d}z,
\end{align*}
其中 $V\triangleq\{(x,y,z)\mid x^2+y^2\leqslant t^2,0\leqslant z\leqslant1\}$.再利用柱坐标换元得
\begin{align*}
I(t)&=\iint_{x^2+y^2\leqslant t^2}\mathrm{d}x\mathrm{d}y\int_0^1 f'\big((x^2+y^2)z\big)\big(2xz+2yz+x^2+y^2\big)\mathrm{d}z\\
&=\int_0^t \mathrm{d}r\int_0^1 \mathrm{d}z\int_0^{2\pi} f'(r^2z)\big(2zr\cos\theta+2zr\sin\theta+r^2\big)\cdot r\mathrm{d}\theta\\
&=\int_0^t \mathrm{d}r\int_0^1 2\pi f'(r^2z)\cdot r^3\mathrm{d}z
=\int_0^t 2\pi r f(r^2)\mathrm{d}r.
\end{align*}
从而再由L'Hospital法则得
\begin{align*}
\lim_{t\to0^+}\frac{I(t)}{t^4}=\lim_{t\to0^+}\frac{I'(t)}{4t^3}=\lim_{t\to0^+}\frac{2\pi t f(t^2)}{4t^3}=\lim_{t\to0^+}\frac{\pi f(t^2)}{2t^2}=\frac{\pi}{2}f'(0).
\end{align*}
\end{solution}

\subsection{利用调和函数性质计算}

\begin{definition}
设$D$为$\mathbb{R}^n$中的开区域，函数$u\in C^n(D)$。若
\begin{align*}
\Delta u\triangleq \sum_{i=1}^n{\frac{\partial ^2u}{\partial x_{i}^{2}}}=0, \quad \forall (x,y)\in D,
\end{align*}
则称$u$为$D$上的\textbf{调和函数},也称$u$在区域$D$上\textbf{调和}.
\end{definition}

\begin{theorem}[调和函数的平均值定理]\label{theorem:数分-调和函数的平均值定理}
设$u$在开区域$D\subset \mathbb{R}^2$内调和，$(x_0,y_0)\in D$，$r>0$使得闭圆盘$\overline{B_r(x_0,y_0)}\subset D$。则
\begin{enumerate}[(1)]
\item 圆周平均值公式：
\begin{align*}
u(x_0,y_0)=\frac{1}{2\pi}\int_{0}^{2\pi} u(x_0+r\cos\theta, y_0+r\sin\theta)\,\mathrm{d}\theta. 
\end{align*}
\item 圆盘平均值公式：
\begin{align*}
u(x_0,y_0)=\frac{1}{\pi r^2}\iint_{B_r(x_0,y_0)} u(x,y)\,\mathrm{d}x\mathrm{d}y. 
\end{align*}
\end{enumerate}
\end{theorem}
\begin{proof}
\begin{enumerate}[(1)]
\item 对$0<\rho<r$，记$B_\rho$为以$(x_0,y_0)$为圆心、$\rho$为半径的圆盘，其边界为$\partial B_\rho$。由\hyperref[theorem:Green第一恒等式]{Green第一恒等式}有
\begin{align*}
0=\iint_{B_\rho} \Delta u\,\mathrm{d}x\mathrm{d}y
=\oint_{\partial B_\rho} \frac{\partial u}{\partial n}\,\mathrm{d}s,
\end{align*}
其中$\frac{\partial u}{\partial n}$为外法向导数。在极坐标$(x,y)=(x_0+\rho\cos\theta, y_0+\rho\sin\theta)$下，$\frac{\partial u}{\partial n}=u_\rho$，$\mathrm{d}s=\rho\,\mathrm{d}\theta$，故由\hyperref[theorem:含参反常积分的可微性]{含参反常积分的可微性}知
\begin{align*}
0=\int_{0}^{2\pi} u_\rho(\rho,\theta)\,\rho\,\mathrm{d}\theta
=\rho \frac{\mathrm{d}}{\mathrm{d}\rho}\left(\int_{0}^{2\pi} u(\rho,\theta)\,\mathrm{d}\theta\right).
\end{align*}
因此
\begin{align*}
\int_{0}^{2\pi} u(\rho,\theta)\,\mathrm{d}\theta = C \in \mathbb{R},
\end{align*}
令$\rho\to 0^+$，由连续性得
\begin{align*}
C=2\pi u(x_0,y_0),
\end{align*}
从而圆周平均值公式成立。

\item 将圆周平均值公式两边乘以$\rho$，并对$\rho$从$0$到$r$积分：
\begin{align*}
u(x_0,y_0)\int_{0}^{r} \rho\,\mathrm{d}\rho
=\frac{1}{2\pi}\int_{0}^{r}\int_{0}^{2\pi} u(x_0+\rho\cos\theta, y_0+\rho\sin\theta)\rho\,\mathrm{d}\theta\mathrm{d}\rho,
\end{align*}
再取$x=x_0+\rho \cos\theta,y=y_0+\rho \sin \theta$换元得
\begin{align*}
u(x_0,y_0)\frac{r^2}{2}
=\frac{1}{2\pi}\iint_{B_r} u\,\mathrm{d}x\mathrm{d}y,
\end{align*}
整理得
\begin{align*}
u(x_0,y_0)=\frac{1}{\pi r^2}\iint_{B_r} u\,\mathrm{d}x\mathrm{d}y.
\end{align*}
从而圆盘平均值公式成立。
\end{enumerate}
\end{proof}


\begin{example}
定义在 $D \triangleq \{(x,y): x^2+y^2 \leqslant R^2\}$ 的函数 $f$ 二阶连续可微, 满足 $f(0,0)=0$, $\Delta f = x + x^2 + y^2$, 计算积分
\begin{align}
I = \iint_D f(x,y)\mathrm{d}x\mathrm{d}y.
\label{eq:double_int_f_disk_20260704}
\end{align}
\end{example}
\begin{note}
可以利用PDE中关于求解Possion方程的方法反解出满足$\Delta P=\Delta f=x^2+y^2+x$的函数$P$.
\end{note}
\begin{solution}
设$r^2=x^2+y^2$，令
\begin{align*}
P(x,y)=\frac{r^4}{16}+\frac{xr^2}{8}=\frac{\left( x^2+y^2 \right) ^2}{16}+\frac{x\left( x^2+y^2 \right)}{8}.
\end{align*}
直接计算得
\begin{align*}
\Delta P = x^2+y^2+x=\Delta f.
\end{align*}
故$h\triangleq f-P$在$D$内调和。

又因$f(0,0)=0$且$P(0,0)=0$，故$h(0,0)=0$。由\hyperref[theorem:数分-调和函数的平均值定理]{调和函数的平均值定理}，对半径为$R$的圆盘$D$，有
\begin{align*}
h(0,0)=\frac{1}{\pi R^2}\iint_D h(x,y)\,\mathrm{d}x\mathrm{d}y.
\end{align*}
因为$h(0,0)=0$，所以
\begin{align*}
\iint_D h\,\mathrm{d}x\mathrm{d}y=0.
\end{align*}
于是
\begin{align*}
I=\iint_D f\,\mathrm{d}x\mathrm{d}y
=\iint_D (P+h)\,\mathrm{d}x\mathrm{d}y
=\iint_D P\,\mathrm{d}x\mathrm{d}y.
\end{align*}
而
\begin{align*}
\iint_D P\,\mathrm{d}x\mathrm{d}y
=\frac{1}{16}\iint_D r^4\,\mathrm{d}x\mathrm{d}y
+\frac{1}{8}\iint_D x r^2\,\mathrm{d}x\mathrm{d}y.
\end{align*}
由于$D$关于$y$轴对称，且$x r^2$关于$x$为奇函数，故
\begin{align*}
\iint_D x r^2\,\mathrm{d}x\mathrm{d}y=0.
\end{align*}
因此
\begin{align*}
I=\frac{1}{16}\iint_D r^4\,\mathrm{d}x\mathrm{d}y.
\end{align*}
用极坐标变换得
\begin{align*}
I=\frac{1}{16}\int_{0}^{2\pi}\int_{0}^{R} \rho^4\cdot \rho\,\mathrm{d}\rho\mathrm{d}\theta 
=\frac{1}{16}\cdot 2\pi\cdot \frac{R^6}{6} =\frac{\pi R^6}{48}.
\end{align*}
\end{solution}














\end{document}